\documentclass[11pt]{article}

\usepackage[margin=1in]{geometry}
\usepackage{xcolor}
\usepackage{amsmath}
\usepackage{graphicx}
\usepackage{listings}
\usepackage{xurl}
\usepackage{hyperref}
\usepackage{booktabs}
\usepackage{enumitem}
\usepackage{parskip}
\usepackage{titlesec}
\usepackage{tcolorbox}
\tcbuselibrary{skins,breakable}

\hypersetup{
  colorlinks = true,
  linkcolor  = blue!60!black,
  urlcolor   = blue!60!black,
  citecolor  = blue!60!black,
}

% --- Code listing style -------------------------------------------------------
\definecolor{shellbg}{HTML}{F5F5F5}
\definecolor{shellrule}{HTML}{CCCCCC}
\definecolor{shellcomment}{HTML}{6A737D}
\definecolor{shellkeyword}{HTML}{005CC5}

\lstdefinestyle{shell}{
  basicstyle      = \ttfamily\footnotesize,
  backgroundcolor = \color{shellbg},
  frame           = single,
  rulecolor       = \color{shellrule},
  framesep        = 6pt,
  xleftmargin     = 6pt,
  xrightmargin    = 6pt,
  breaklines      = true,
  showstringspaces= false,
  columns         = fullflexible,
  keepspaces      = true,
  commentstyle    = \color{shellcomment}\itshape,
  keywordstyle    = \color{shellkeyword}\bfseries,
  morecomment     = [l]{\#},
  morekeywords    = {xschem, ngspice, magic, netgen, klayout, export, source,
                     cd, git, sudo, apt-get, dnf, tran, dc, ac, plot, run},
  literate        = {/}{{/}}1 {-}{{-}}1 {=}{{=}}1,
}
\lstset{style=shell}

% --- Tip / Warning boxes ------------------------------------------------------
\newtcolorbox{tipbox}{
  colback  = blue!5!white,
  colframe = blue!60!black,
  fonttitle= \bfseries,
  title    = Tip,
  breakable,
}
\newtcolorbox{warnbox}{
  colback  = orange!8!white,
  colframe = orange!80!black,
  fonttitle= \bfseries,
  title    = Important,
  breakable,
}

\titleformat{\section}{\Large\bfseries}{\thesection}{0.7em}{}
\titleformat{\subsection}{\large\bfseries}{\thesubsection}{0.6em}{}
\titleformat{\subsubsection}{\normalsize\bfseries}{\thesubsubsection}{0.5em}{}

% ==============================================================================
\title{A Practical Tutorial for Xschem Schematic Capture\\[4pt]
       \large Using the GlobalFoundries 180\,nm PDK (\texttt{gf180mcuD})}
\author{James E.\ Stine and Landon Burleson\\
        Oklahoma State University\\
        Electrical and Computer Engineering Department\\
        Stillwater, OK 74078 USA\\
        \texttt{james.stine@okstate.edu}}
\date{\today}
% ==============================================================================

\begin{document}
\maketitle
\tableofcontents
\newpage

% ==============================================================================
\section{Introduction}
% ==============================================================================

Xschem is an open-source schematic capture and netlist generation tool
originally developed by Stefan Schippers and now maintained at
\url{https://xschem.sourceforge.io}.
It is the de facto schematic editor for the open-source analog/mixed-signal
IC design flow, providing fast hierarchical schematic entry, parameterized
symbol creation, and direct integration with ngspice, Xyce, and other
SPICE-class simulators.
Xschem receives frequent updates; always prefer version 3.4 or later.

This tutorial targets the \textbf{GlobalFoundries 180\,nm MCU PDK}
(\texttt{gf180mcuD}), an open-source process development kit available
through the \texttt{open\_pdks} installer and supported by Google's
open-silicon initiative.
The \texttt{gf180mcuD} library provides Xschem-ready symbols for nMOS,
pMOS, resistors, capacitors, BJTs, and diodes, along with model files for
ngspice simulation.
All symbol names, parameter conventions, and command examples in this
document refer to \texttt{gf180mcuD} unless otherwise noted.

\begin{tipbox}
Always sketch your schematic on paper or a whiteboard before opening
Xschem.
A clean schematic is a planning artifact, not a drawing exercise:
spending five minutes thinking through power and ground placement,
signal flow, and naming conventions will save you hours of rewiring,
deleting, and renaming inside the editor.
A good rule is \emph{signals flow left-to-right, supplies run top
(VDD) to bottom (GND), and inputs land on the left edge while outputs
exit on the right.}
Once that mental layout is fixed, schematic entry becomes mechanical.
\end{tipbox}

% ==============================================================================
\section{Prerequisites and Setup}
% ==============================================================================

\subsection{Software}

You need Xschem (3.4+), ngspice (40+), and \texttt{open\_pdks} with
\texttt{gf180mcuD} installed.
If you are using the IIC-OSIC-TOOLS Docker container
(\url{https://github.com/iic-jku/iic-osic-tools}), all of this is
pre-installed and pre-configured---skip ahead to
Section~\ref{sec:launching-xschem}.

For a native Linux build, follow the accompanying installation guide
and confirm that the following environment variables are set
(\texttt{setup.sh} handles this automatically):

\begin{lstlisting}
export OPEN_CIRCUITDESIGN_ROOT="$HOME/ocd"
export PDK_ROOT="$OPEN_CIRCUITDESIGN_ROOT/share/pdk"
export PDK=gf180mcuD
\end{lstlisting}

\subsection{Environment Check}

Before starting any design work, verify that Xschem can see the PDK
symbols:

\begin{lstlisting}
ls $PDK_ROOT/gf180mcuD/libs.tech/xschem/
# Should list: gf180mcu, sources, simulator_commands, etc.
\end{lstlisting}

If the directory is empty or missing, re-run \texttt{open\_pdks} for
\texttt{gf180mcuD} or verify that \texttt{PDK\_ROOT} is set correctly.

\subsection{Configuring \texttt{xschemrc}}

Xschem reads its configuration from an \texttt{xschemrc} file, which
tells the tool where to find PDK symbol libraries, simulation models,
and your personal designs.
The PDK ships with one ready to use:

\begin{lstlisting}
ls $PDK_ROOT/gf180mcuD/libs.tech/xschem/xschemrc
\end{lstlisting}

You can either invoke Xschem with that file directly or copy it to
\texttt{\textasciitilde/.xschem/xschemrc} and customize the
\texttt{XSCHEM\_LIBRARY\_PATH} variable to include your own design
directories.

\subsection{Create a Working Directory}

Organize each design in its own directory to avoid file collisions:

\begin{lstlisting}
mkdir -p ~/designs/inverter
cd ~/designs/inverter
\end{lstlisting}

% ==============================================================================
\section{Launching Xschem}
\label{sec:launching-xschem}
% ==============================================================================

Always launch Xschem with the \texttt{gf180mcuD} configuration so that
the correct symbol libraries and simulation models are loaded:

\begin{lstlisting}
xschem --rcfile $PDK_ROOT/gf180mcuD/libs.tech/xschem/xschemrc
\end{lstlisting}

Xschem opens two windows simultaneously:

\begin{itemize}[leftmargin=2em]
  \item \textbf{Schematic (drawing) window} --- where you place
        components and draw wires.
  \item \textbf{Terminal (console) window} --- the shell where Xschem
        was launched, which displays netlist generation messages,
        simulator output, and all error/warning text.
\end{itemize}

\begin{warnbox}
Keep \emph{both} windows visible at all times.
Symbol-load failures, parameter parsing errors, and simulator messages
appear \emph{only} in the terminal window.
The schematic window will look completely normal even when serious
problems exist---a missing model file, a malformed parameter, or a
simulation that exited immediately leaves no visible artifact in the
schematic itself.
The \emph{only} indication that something has gone wrong is a message
in the terminal.
A few practical habits follow from this:
\begin{itemize}
\item \textbf{After every netlist generation} (\texttt{Netlist} menu or
      pressing \texttt{N}), glance at the terminal before continuing.
\item \textbf{After launching a simulation}, scroll through the full
      ngspice output rather than stopping at the first ``OK'' line.
\item \textbf{When behavior seems wrong}---a component appears blank, a
      simulation produces no waveforms, or a launch button does
      nothing---the explanation is almost always already waiting in the
      terminal window.
\end{itemize}
Treat the terminal not as a log to be consulted occasionally, but as a
live diagnostic stream that runs in parallel with everything you do.
\end{warnbox}

% ==============================================================================
\section{The Xschem Interface}
% ==============================================================================

\subsection{Mouse Buttons}

The mouse is your primary schematic instrument.
Learn these three actions before anything else:

\begin{center}
\begin{tabular}{@{}llp{8cm}@{}}
\toprule
\textbf{Button} & \textbf{Shorthand} & \textbf{Action} \\
\midrule
Left click   & L & Place an instance, end a wire, or select an item \\
Middle click & M & Pan the view; drag with middle-button to pan freely \\
Right click  & R & Open the context menu (properties, delete, copy, etc.) \\
\bottomrule
\end{tabular}
\end{center}

The mouse wheel zooms in and out about the cursor location.
A drag with left-click on empty canvas creates a rubber-band selection
rectangle---everything fully enclosed becomes selected.

\begin{tipbox}
Get used to right-clicking on instances and wires.
The context menu is the fastest path to editing properties, copying
selections, and deleting items, and it always shows you only the
operations that are valid for the object under the cursor.
\end{tipbox}

\subsection{Entering Commands}

Xschem combines several input modes:

\begin{enumerate}[leftmargin=2em]
  \item \textbf{Single-character keyboard macros} typed in the
        schematic window (e.g., \texttt{w} starts a wire,
        \texttt{Insert} opens the symbol browser, \texttt{f} fits the
        view).
  \item \textbf{Menu commands} from the top menu bar (File, Edit, View,
        Symbol, Simulation, etc.).
  \item \textbf{Mouse actions} for placement, selection, and panning.
\end{enumerate}

The most important macros to learn first are \texttt{Insert} (place a
new instance), \texttt{w} (wire), \texttt{l} (label/wire name),
\texttt{Escape} (cancel current operation), and \texttt{Ctrl+s} (save).

\subsection{Selection and Movement}

Click to select a single object.
Hold \texttt{Shift} and click to add to the selection, or drag a
rubber-band rectangle to select a region.
Move selected items with \texttt{m}, copy them with \texttt{c}, and
delete them with \texttt{Delete}.
Press \texttt{Escape} at any time to cancel the current operation
without committing changes.
\clearpage

\begin{warnbox}
A common source of confusion is starting a wire with \texttt{w} and
then accidentally beginning to type a label or instance name into the
schematic window.
Xschem captures every keystroke as a potential macro, so what you
\emph{think} is text entry can instead trigger another command---moving,
deleting, or rotating the selection unexpectedly.
If something stops working as expected---a click produces no result, or
the cursor seems stuck in some other mode---press \texttt{Escape} once
or twice to return to the default state before doing anything else.
Make it a habit to press \texttt{Escape} between distinct operations,
particularly after wire drawing, label placement, or property editing.
\end{warnbox}

% ==============================================================================
\section{Essential Command Reference}
% ==============================================================================

\subsection{File and View Commands}

\begin{center}
\begin{tabular}{@{}lp{9.5cm}@{}}
\toprule
\textbf{Command/Key} & \textbf{Description} \\
\midrule
\texttt{Ctrl+s}                & Save current schematic \\
\texttt{Ctrl+o}                & Open existing schematic \\
\texttt{Ctrl+n}                & New empty schematic in a fresh window \\
\texttt{Ctrl+q}                & Quit Xschem \\
\texttt{f}                     & Fit all geometry to window (zoom-fit) \\
\texttt{z} / \texttt{Z}        & Zoom in / zoom out about the cursor \\
\texttt{Ctrl+e}                & Edit (descend into) the selected
                                 instance's schematic \\
\texttt{Ctrl+i}                & Edit (descend into) the selected
                                 instance's symbol \\
\texttt{Ctrl+a}                & Select all \\
\texttt{Esc}                   & Cancel current operation \\
\bottomrule
\end{tabular}
\end{center}

\subsection{Drawing and Editing}

\begin{center}
\begin{tabular}{@{}lll@{}}
\toprule
\textbf{Action}            & \textbf{Macro}    & \textbf{Description} \\
\midrule
Place new instance         & \texttt{Insert}   & Open symbol browser to
                                                 place a component \\
Draw wire                  & \texttt{w}        & Start a wire; left-click
                                                 to add vertices \\
Place label/wire name      & \texttt{l}        & Insert a net label
                                                 (\texttt{lab\_pin}) \\
Place text annotation      & \texttt{t}        & Add free-form text
                                                 (non-electrical) \\
Move selection             & \texttt{m}        & Move selected items \\
Copy selection             & \texttt{c}        & Copy selected items \\
Rotate selection           & \texttt{r}        & Rotate 90$^\circ$
                                                 counter-clockwise \\
Flip selection             & \texttt{Shift+f}  & Mirror selection \\
Delete selection           & \texttt{Delete}   & Remove selected items \\
Edit properties            & \texttt{q}        & Open property dialog
                                                 for selected instance \\
\bottomrule
\end{tabular}
\end{center}

\subsection{Netlist and Simulation}

\begin{center}
\begin{tabular}{@{}lp{9.5cm}@{}}
\toprule
\textbf{Macro/Menu}         & \textbf{Description} \\
\midrule
\texttt{Netlist}            & Generate a SPICE netlist for the current
                              schematic \\
\texttt{Simulate}           & Launch the configured simulator (ngspice
                              by default) on the netlist \\
\texttt{a}                  & Show/hide net annotations after a
                              simulation \\
\texttt{Edit netlist}       & Open the most recent netlist in your
                              editor for inspection \\
\bottomrule
\end{tabular}
\end{center}

The \texttt{Options} menu lets you switch netlist format between SPICE,
Verilog, VHDL, and tEDAx.
For \texttt{gf180mcuD} simulation work, leave it on SPICE.

% ==============================================================================
\section{gf180mcuD Symbol Reference}
% ==============================================================================

The PDK provides Xschem symbols for every primitive device.
The most commonly used symbols for digital and mixed-signal cell design
are listed below; the \emph{Symbol path} is what appears in the
\texttt{Insert} symbol browser.

\begin{center}
\begin{tabular}{@{}lp{9cm}@{}}
\toprule
\textbf{Symbol Path}                          & \textbf{Description} \\
\midrule
\texttt{gf180mcu/nfet\_03v3}                  & Standard 3.3\,V nMOS
                                                transistor \\
\texttt{gf180mcu/pfet\_03v3}                  & Standard 3.3\,V pMOS
                                                transistor \\
\texttt{gf180mcu/nfet\_06v0}                  & 5.0\,V (6\,V tolerant)
                                                nMOS \\
\texttt{gf180mcu/pfet\_06v0}                  & 5.0\,V (6\,V tolerant)
                                                pMOS \\
\texttt{gf180mcu/res\_*}                      & Resistor primitives
                                                (poly, metal, diffusion) \\
\texttt{gf180mcu/cap\_*}                      & Capacitor primitives
                                                (MIM, MOS, MOM) \\
\texttt{gf180mcu/diode\_*}                    & Diode primitives \\
\texttt{devices/ipin}                         & Schematic input port
                                                (left edge of cell) \\
\texttt{devices/opin}                         & Schematic output port
                                                (right edge of cell) \\
\texttt{devices/iopin}                        & Bidirectional port \\
\texttt{devices/lab\_pin}                     & Net-name label \\
\texttt{devices/gnd}                          & Ground symbol \\
\texttt{devices/vsource}                      & Voltage source \\
\texttt{devices/code\_shown}                  & Inline SPICE code block
                                                (for \texttt{.tran},
                                                \texttt{.include}, etc.) \\
\bottomrule
\end{tabular}
\end{center}

\begin{tipbox}
\textbf{Use the right transistor flavor for your supply.}
\texttt{gf180mcuD} provides separate device families for 3.3\,V and
5\,V operation.
Mixing them in a single design is allowed---and often necessary for
mixed-signal interfaces---but each device must operate within its rated
\(V_{ds}\), \(V_{gs}\), and \(V_{bs}\) limits or it will degrade quickly
in silicon.
For most digital cell work in \texttt{gf180mcuD}, the 3.3\,V family
(\texttt{nfet\_03v3} and \texttt{pfet\_03v3}) is the correct choice.
Reach for the 5\,V devices only when an interface specifically requires
them.
\end{tipbox}

% ==============================================================================
\section{Ports, Labels, and Hierarchy}
% ==============================================================================

\subsection{Ports vs.\ Labels}

Xschem distinguishes between two kinds of named connections:

\begin{itemize}[leftmargin=2em]
  \item \textbf{Ports} (\texttt{ipin}, \texttt{opin}, \texttt{iopin})
        define the external interface of a cell.
        When you generate a symbol from this schematic, every port
        becomes a pin on the symbol.
  \item \textbf{Labels} (\texttt{lab\_pin}) name internal nets so they
        can be referred to by name elsewhere in the same schematic, or
        so they appear meaningfully in the netlist.
\end{itemize}

A single net can carry multiple labels with the same name; Xschem
treats them as electrically connected even without a visible wire.
This is the standard way to wire up power and ground rails without
cluttering the schematic.

\subsection{Naming Conventions}

\begin{itemize}[leftmargin=2em]
  \item \textbf{\texttt{VDD}} and \textbf{\texttt{VSS}}
        (or \textbf{\texttt{GND}}) --- supply rails.
        Use the same names as in your testbench and any other
        cells you intend to instantiate.
  \item Signal names must match between the schematic and its companion
        layout for LVS to pass; use exactly the same capitalization.
  \item Adopt a consistent style within a project---\texttt{a},
        \texttt{b}, \texttt{y}, \texttt{clk}, \texttt{rst}, etc.---and
        keep it everywhere.
\end{itemize}

\begin{tipbox}
\textbf{Match labels to the schematic of the surrounding system.}
A net labeled \texttt{Y} in your inverter must connect to a net also
labeled \texttt{Y} in any cell that uses it.
LVS, top-level netlist generation, and inter-cell wiring all depend on
character-for-character label matches.
A net labeled \texttt{vdd} will not connect to a net labeled
\texttt{VDD}, and Xschem will not warn you about the mismatch---it will
simply produce a netlist with two disconnected supplies.
\end{tipbox}

\subsection{Creating a Symbol from a Schematic}

Every schematic with ports can be turned into a symbol that other
schematics can instantiate.
With your schematic open, choose \textbf{Symbol $\rightarrow$ Make
symbol from schematic} (or press \texttt{a} from the menu shortcut).
Xschem auto-generates a rectangular symbol with one pin per port,
saved as \texttt{<name>.sym} alongside your \texttt{<name>.sch}.
After that, instantiate it from another schematic with \texttt{Insert}
and browse to your working directory.

\subsection{Descending and Ascending the Hierarchy}

Select an instance and press \texttt{Ctrl+e} to descend into its
schematic, or \texttt{Ctrl+i} to descend into its symbol drawing.
Press \texttt{Ctrl+e} again on empty canvas (or use the \texttt{<<}
button on the toolbar) to ascend back to the parent.

% ==============================================================================
\section{Step-by-Step: CMOS Inverter Schematic in gf180mcuD}
% ==============================================================================

This section walks through drawing a minimum-size CMOS inverter from
scratch in \texttt{gf180mcuD}.
Have your sketch ready before you begin.

\subsection{Transistor Sizing}

In \texttt{gf180mcuD} the minimum drawn gate length is
$L_{\min} = 0.18\,\mu\text{m}$.
A reasonable starting inverter uses:

\begin{center}
\begin{tabular}{@{}llll@{}}
\toprule
\textbf{Transistor} & \textbf{Type} & \textbf{Width} & \textbf{Length} \\
\midrule
nMOS & \texttt{nfet\_03v3} & $0.42\,\mu\text{m}$ & $0.18\,\mu\text{m}$ \\
pMOS & \texttt{pfet\_03v3} & $0.84\,\mu\text{m}$ ($2\times$ nMOS) & $0.18\,\mu\text{m}$ \\
\bottomrule
\end{tabular}
\end{center}

The pMOS is sized $2\times$ wider than the nMOS to equalize rise and
fall times, since hole mobility is roughly half that of electrons.

\subsection{Step 1 --- Launch Xschem and Create a New Schematic}

\begin{lstlisting}
cd ~/designs/inverter
xschem --rcfile $PDK_ROOT/gf180mcuD/libs.tech/xschem/xschemrc
\end{lstlisting}

In the schematic window, press \texttt{Ctrl+n} to open a new empty
schematic.
Save it immediately so all subsequent operations have a target file:

\begin{lstlisting}
File -> Save As ...   # save as inv.sch in ~/designs/inverter
\end{lstlisting}

\subsection{Step 2 --- Place the nMOS Transistor}

Press \texttt{Insert} to open the symbol browser, then navigate to
\texttt{gf180mcu/nfet\_03v3} and double-click to attach it to the
cursor.
Left-click in the lower half of the canvas to place it.
Press \texttt{Escape} to exit placement mode.

Select the placed instance, press \texttt{q} to open its properties,
and set:

\begin{lstlisting}
W = 0.42u
L = 0.18u
nf = 1
\end{lstlisting}

\subsection{Step 3 --- Place the pMOS Transistor}

Press \texttt{Insert} again, navigate to \texttt{gf180mcu/pfet\_03v3},
and place it directly above the nMOS.
Edit its properties (\texttt{q}):

\begin{lstlisting}
W = 0.84u
L = 0.18u
nf = 1
\end{lstlisting}

\subsection{Step 4 --- Place Power, Ground, and Substrate Connections}

A MOSFET symbol in the PDK has four terminals: drain, gate, source, and
body.
The body terminals must be tied to the appropriate supply for correct
operation:

\begin{itemize}[leftmargin=2em]
  \item nMOS body $\rightarrow$ \texttt{VSS} (ground).
  \item pMOS body (n-well) $\rightarrow$ \texttt{VDD}.
\end{itemize}

The simplest way to make these connections is to label the body wires
with \texttt{VDD} and \texttt{VSS} so they connect by name instead of
running explicit wires.

\subsection{Step 5 --- Draw the Wires}

Press \texttt{w} to enter wire mode and connect:

\begin{itemize}[leftmargin=2em]
  \item pMOS source $\rightarrow$ \texttt{VDD} rail.
  \item nMOS source $\rightarrow$ \texttt{VSS} rail.
  \item pMOS drain + nMOS drain $\rightarrow$ output node.
  \item pMOS gate + nMOS gate $\rightarrow$ input node.
\end{itemize}

Click each vertex of the wire path; double-click or press \texttt{Esc}
to finish.
Wires snap to the grid automatically and form T-junctions wherever they
cross other wires at vertices.

\begin{tipbox}
\textbf{Always wire on grid.}
Xschem will not connect two wires that pass each other off-grid,
even if they look connected on screen.
If a node refuses to take a label or shows up disconnected after
netlisting, the most common cause is one endpoint sitting half a grid
unit off.
Press \texttt{f} to fit, then zoom in (\texttt{z}) on the suspected
junction and verify both endpoints are exactly on grid intersections.
\end{tipbox}

\subsection{Step 6 --- Add Port Symbols}

Press \texttt{Insert}, browse to \texttt{devices/ipin}, and place an
input pin on the input wire.
Edit its properties (\texttt{q}) and set:

\begin{lstlisting}
name = A
\end{lstlisting}

Repeat with \texttt{devices/opin} on the output wire, naming it
\texttt{Y}.
These ports define the external interface of the cell---when you
generate a symbol later, \texttt{A} and \texttt{Y} become the input
and output pins.

\subsection{Step 7 --- Add Labels for Power and Ground}

Press \texttt{Insert}, browse to \texttt{devices/lab\_pin}, and place
labels on the supply wires:

\begin{lstlisting}
name = VDD       # on the upper rail and pMOS body
name = VSS       # on the lower rail and nMOS body
\end{lstlisting}

Anywhere two \texttt{lab\_pin} instances share the same name, Xschem
treats the wires they sit on as electrically connected---even with no
visible wire between them.

\subsection{Step 8 --- Save}

\begin{lstlisting}
Ctrl+s
\end{lstlisting}

This writes \texttt{inv.sch} to your working directory.
Save frequently; Xschem does not auto-save.

\subsection{Step 9 --- Generate the Symbol}

With the schematic saved, choose \textbf{Symbol $\rightarrow$ Make
symbol from schematic}.
Xschem creates \texttt{inv.sym} in the same directory, ready to be
instantiated by a testbench or higher-level cell.

\begin{warnbox}
If you later add, remove, or rename ports, regenerate the symbol.
A schematic and its symbol can drift out of sync silently---an
\texttt{ipin} added to the schematic will not appear on the symbol
until you regenerate it, and any cell that already instantiated the old
symbol will continue to use the old pin list with no warning.
Make symbol regeneration part of your save workflow whenever the port
list changes.
\end{warnbox}

% ==============================================================================
\section{Building a Testbench}
% ==============================================================================

A testbench schematic instantiates the cell under test, applies stimulus,
and includes the simulator directives needed to run a particular
analysis.
This section builds a minimal inverter testbench.

\subsection{Step 1 --- New Schematic}

Start a fresh schematic with \texttt{Ctrl+n} and save it as
\texttt{inv\_tb.sch} in the same directory as \texttt{inv.sch}.

\subsection{Step 2 --- Instantiate the Inverter}

Press \texttt{Insert} and browse to your working directory rather than
the PDK library.
Select \texttt{inv.sym} and place it in the center of the canvas.

\subsection{Step 3 --- Add the Supply}

Place a \texttt{devices/vsource} symbol, set its properties to
\texttt{value = 3.3} and \texttt{name = VDD}, and wire its positive
terminal to a label \texttt{VDD} and its negative terminal to a label
\texttt{VSS}.
Place a \texttt{devices/gnd} symbol on the \texttt{VSS} node so that
ngspice has a defined zero-volt reference.

\subsection{Step 4 --- Add the Input Stimulus}

Place another \texttt{devices/vsource} on input \texttt{A}.
Edit its properties:

\begin{lstlisting}
name  = VIN
value = pulse(0 3.3 0 100p 100p 5n 10n)
\end{lstlisting}

This drives \texttt{A} with a 100\,MHz pulse train (10\,ns period,
50\% duty cycle, 100\,ps edges).
The negative terminal connects to \texttt{VSS}.

\subsection{Step 5 --- Add Simulator Directives}

Place a \texttt{devices/code\_shown} symbol anywhere on the canvas.
Edit its \texttt{value} property to contain the simulation control
deck:

\begin{lstlisting}
.lib "$PDK_ROOT/gf180mcuD/libs.tech/ngspice/design.ngspice" typical
.tran 10p 50n
.save all
.control
run
plot v(A) v(Y)
.endc
\end{lstlisting}

The first line includes the typical-corner device models for
\texttt{gf180mcuD}.
The \texttt{.tran} directive runs a 50\,ns transient with 10\,ps
internal time steps, the \texttt{.save all} stores every node, and the
\texttt{.control} block automatically plots the input and output once
the simulation finishes.

\subsection{Step 6 --- Generate the Netlist}

Press \texttt{N} (or use \textbf{Simulation $\rightarrow$ Netlist}).
Xschem writes \texttt{inv\_tb.spice} to the working directory and
prints any errors to the terminal.

\subsection{Step 7 --- Run the Simulation}

Press the \textbf{Simulate} button (or use the \textbf{Simulation}
menu).
Xschem launches ngspice on the netlist; ngspice runs in a console and
opens a plot window with the requested waveforms.

\begin{tipbox}
\textbf{Re-netlist before every simulation.}
Xschem does not automatically regenerate the netlist when you change
the schematic.
A common debugging trap is to edit the schematic, hit Simulate, and see
exactly the same (broken) result as before---because ngspice ran the
old netlist.
Make \texttt{N} (netlist) a reflex before \texttt{Simulate}, or use the
\textbf{Netlist + Simulate} menu option that does both in one step.
\end{tipbox}

% ==============================================================================
\section{Netlist Generation in Detail}
% ==============================================================================

The netlist Xschem generates is a plain-text SPICE file usable directly
by ngspice, Xyce, or any other Berkeley-compatible simulator.
You can inspect it with any text editor:

\begin{lstlisting}
cat inv_tb.spice
\end{lstlisting}

Useful options under the \textbf{Options} menu:

\begin{itemize}[leftmargin=2em]
  \item \textbf{Spice netlist}: standard SPICE output (default).
  \item \textbf{Hierarchical netlist}: emits subcircuit definitions
        (\texttt{.subckt}) for every cell rather than flattening.
        Required for any design with more than one level of hierarchy.
  \item \textbf{Top-level is subckt}: wraps the top schematic itself
        in a \texttt{.subckt} block so it can be reused.
        Leave this off for testbenches.
\end{itemize}

\begin{warnbox}
For LVS-clean cell design, always use the hierarchical netlister.
A flattened netlist destroys the cell-by-cell structure that LVS tools
rely on to compare your schematic against the layout.
Magic's \texttt{.spice} extraction is hierarchical by default; your
schematic netlist must match.
\end{warnbox}

% ==============================================================================
\section{Simulation with ngspice}
% ==============================================================================

When Xschem launches ngspice, you get an interactive console where you
can issue additional commands after the initial \texttt{.control} block
finishes:

\begin{lstlisting}
ngspice 1 -> plot v(A) v(Y)
ngspice 2 -> meas tran t_pdr trig v(A) val=1.65 fall=1 \
                                targ v(Y) val=1.65 rise=1
ngspice 3 -> meas tran t_pdf trig v(A) val=1.65 rise=1 \
                                targ v(Y) val=1.65 fall=1
ngspice 4 -> let t_pd = (t_pdr + t_pdf) / 2
ngspice 5 -> print t_pd
\end{lstlisting}

These four \texttt{meas} statements compute the rising-edge propagation
delay (\texttt{t\_pdr}), falling-edge propagation delay
(\texttt{t\_pdf}), and their average (\texttt{t\_pd}) for the inverter.

\subsection{Common Analyses}

\begin{center}
\begin{tabular}{@{}lp{8.5cm}@{}}
\toprule
\textbf{Directive} & \textbf{Purpose} \\
\midrule
\texttt{.tran 10p 50n}        & Transient: 10\,ps step, 50\,ns stop \\
\texttt{.dc VIN 0 3.3 0.01}   & DC sweep of \texttt{VIN} from 0 to
                                3.3\,V in 10\,mV steps \\
\texttt{.ac dec 10 1 1G}      & AC sweep, 10 points per decade,
                                1\,Hz to 1\,GHz \\
\texttt{.noise v(out) VIN dec 10 1 1G}
                              & Noise analysis at output referred to
                                \texttt{VIN} \\
\texttt{.op}                  & DC operating point only \\
\bottomrule
\end{tabular}
\end{center}

For a simple voltage-transfer-curve plot of an inverter, use
\texttt{.dc} sweeping the input source rather than \texttt{.tran}.
Pulse stimuli can be useful but are unnecessary for static
characterization.

% ==============================================================================
\section{Schematic vs.\ Layout (LVS Workflow)}
% ==============================================================================

Once your schematic simulates correctly, the SPICE netlist Xschem
generates becomes the reference for Layout vs.\ Schematic (LVS)
comparison against the corresponding Magic layout.
Many professional organizations list skipping LVS as a termination
offense.

A typical workflow:

\begin{enumerate}[leftmargin=2em]
  \item In Xschem, generate a hierarchical netlist of the cell
        (\emph{not} the testbench): \texttt{Options $\rightarrow$
        Hierarchical netlist}, then \texttt{Netlist}.
  \item In Magic, extract the layout to SPICE
        (\texttt{:extract all; :ext2spice}).
  \item Run Netgen with the PDK setup file:
\begin{lstlisting}
netgen -batch lvs "inv.spice inv" \
    "$PDK_ROOT/gf180mcuD/libs.tech/netgen/gf180mcuD_setup.tcl"
\end{lstlisting}
  \item A passing comparison prints \textbf{Circuits match}.
        A failure lists mismatched devices and nets---trace each one
        back to a label or wiring error in either side.
\end{enumerate}

\begin{tipbox}
The schematic is the \emph{golden} reference.
If LVS fails, the layout is what changes---never edit the schematic to
``match'' a faulty layout, since the schematic is what your simulation
verified.
The only exception is when the schematic itself is wrong relative to
the design intent, in which case you fix the schematic, re-simulate,
re-netlist, and re-run LVS from the top.
\end{tipbox}

% ==============================================================================
\section{Tips and Common Pitfalls}
% ==============================================================================

\begin{itemize}[leftmargin=2em]
  \item \textbf{Re-netlist before every simulation.}
        Xschem does not automatically regenerate the netlist when you
        edit the schematic.
        Press \texttt{N} before \texttt{Simulate} every time, or use
        the combined menu option.

  \item \textbf{Match labels exactly.}
        \texttt{VDD} and \texttt{vdd} are different nets.
        Pick a convention (typically all-uppercase for supplies and
        lowercase for signals) and follow it everywhere.

  \item \textbf{Tie every body terminal.}
        A floating MOSFET body produces nonsense simulation results
        and an LVS failure.
        Always label the body net to \texttt{VDD} (pMOS) or \texttt{VSS}
        (nMOS) explicitly.

  \item \textbf{Keep wires on grid.}
        Off-grid wires look connected on screen but produce
        disconnected nets in the netlist.
        When in doubt, zoom in and verify each junction.

  \item \textbf{Regenerate the symbol after port changes.}
        Adding, removing, or renaming \texttt{ipin}/\texttt{opin}
        instances does not automatically update the symbol file.
        Use \textbf{Symbol $\rightarrow$ Make symbol from schematic}
        whenever the port list changes.

  \item \textbf{Use \texttt{code\_shown} for SPICE directives.}
        Inline SPICE blocks (\texttt{.tran}, \texttt{.lib},
        \texttt{.control}) belong in a \texttt{code\_shown} symbol on
        the testbench schematic, not buried in the netlist.

  \item \textbf{Save your testbench separately.}
        Keep \texttt{<cell>.sch} and \texttt{<cell>\_tb.sch} as
        independent files in the same directory.
        The cell schematic stays clean and reusable; the testbench
        carries the simulation setup.

  \item \textbf{Read the terminal.}
        Almost every diagnostic Xschem and ngspice produce appears in
        the launching terminal, not in the GUI.
        Make a habit of glancing at it after every netlist and
        simulate.

  \item \textbf{Use scripts.}
        Repeated simulation runs, sweep generation, and corner-by-corner
        netlisting are all scriptable from the Tcl console (open with
        \textbf{Window $\rightarrow$ Tcl prompt}).
        A small library of reusable Tcl snippets pays for itself
        quickly.
\end{itemize}

% ==============================================================================
\section{Quick-Reference Cheat Sheet}
% ==============================================================================

\subsection*{Keyboard Macros}

\begin{center}
\begin{tabular}{@{}lll@{}}
\toprule
\textbf{Key} & \textbf{Action} & \textbf{Notes} \\
\midrule
\texttt{Insert}  & Place new instance     & Opens symbol browser \\
\texttt{w}       & Draw wire              & Click vertices; \texttt{Esc} to end \\
\texttt{l}       & Place label            & Insert \texttt{lab\_pin} \\
\texttt{t}       & Place text             & Non-electrical annotation \\
\texttt{m}       & Move selection         & --- \\
\texttt{c}       & Copy selection         & --- \\
\texttt{r}       & Rotate $90^\circ$ CCW  & --- \\
\texttt{Shift+f} & Flip                   & Mirror selection \\
\texttt{q}       & Edit properties        & Of selected instance \\
\texttt{Delete}  & Delete selection       & --- \\
\texttt{f}       & Fit to window          & --- \\
\texttt{z}/\texttt{Z} & Zoom in/out       & About cursor \\
\texttt{Ctrl+e}  & Descend into schematic & Of selected instance \\
\texttt{Ctrl+i}  & Descend into symbol    & Of selected instance \\
\texttt{Ctrl+s}  & Save                   & --- \\
\texttt{Ctrl+a}  & Select all             & --- \\
\texttt{Esc}     & Cancel                 & Always safe \\
\bottomrule
\end{tabular}
\end{center}

\subsection*{Key Commands}

\noindent To launch Xschem with the gf180mcuD configuration:
\begin{lstlisting}
xschem --rcfile $PDK_ROOT/gf180mcuD/libs.tech/xschem/xschemrc
\end{lstlisting}

\noindent Alternatively, copy the PDK \texttt{xschemrc} to
\texttt{\textasciitilde/.xschem/xschemrc} so Xschem loads it
automatically every time:
\begin{lstlisting}
mkdir -p ~/.xschem
cp $PDK_ROOT/gf180mcuD/libs.tech/xschem/xschemrc ~/.xschem/xschemrc
\end{lstlisting}
\noindent With this in place, simply typing \texttt{xschem} from any
directory will load the \texttt{gf180mcuD} symbol libraries without
requiring the \texttt{--rcfile} argument each time.

\begin{center}
\begin{tabular}{@{}lp{7.5cm}@{}}
\toprule
\textbf{Menu/Action} & \textbf{Purpose} \\
\midrule
\texttt{Netlist}                    & Generate SPICE netlist for current
                                      schematic \\
\texttt{Simulate}                   & Run ngspice on the netlist \\
\texttt{Symbol $\to$ Make symbol}   & Generate \texttt{.sym} from
                                      \texttt{.sch} \\
\texttt{Options $\to$ Hierarchical} & Toggle hierarchical netlister
                                      (required for LVS) \\
\texttt{Window $\to$ Tcl prompt}    & Open the Tcl scripting console \\
\bottomrule
\end{tabular}
\end{center}

\subsection*{Common Symbols for gf180mcuD}

\begin{center}
\begin{tabular}{@{}ll@{}}
\toprule
\textbf{Symbol path} & \textbf{What it places} \\
\midrule
\texttt{gf180mcu/nfet\_03v3}   & 3.3\,V nMOS \\
\texttt{gf180mcu/pfet\_03v3}   & 3.3\,V pMOS \\
\texttt{gf180mcu/res\_*}       & Resistor primitives \\
\texttt{gf180mcu/cap\_*}       & Capacitor primitives \\
\texttt{devices/ipin}          & Input port \\
\texttt{devices/opin}          & Output port \\
\texttt{devices/iopin}         & Bidirectional port \\
\texttt{devices/lab\_pin}      & Net label \\
\texttt{devices/gnd}           & Ground \\
\texttt{devices/vsource}       & Voltage source \\
\texttt{devices/code\_shown}   & SPICE directive block \\
\bottomrule
\end{tabular}
\end{center}

% ==============================================================================
\section*{Additional Resources}
\addcontentsline{toc}{section}{Additional Resources}
% ==============================================================================

\begin{itemize}[leftmargin=2em]
  \item Xschem manual and tutorials:
        \url{https://xschem.sourceforge.io/stefan/index.html}
  \item ngspice manual:
        \url{https://ngspice.sourceforge.io/docs.html}
  \item gf180mcuD PDK documentation:
        \url{https://gf180mcu-pdk.readthedocs.io}
  \item open\_pdks repository (PDK installer):
        \url{https://github.com/RTimothyEdwards/open_pdks}
  \item IIC-OSIC-TOOLS Docker container:
        \url{https://github.com/iic-jku/iic-osic-tools}
  \item Course examples, scripts, and schematic templates:
        \url{https://stineje.github.io}
  \item Video walkthroughs:
        \url{https://www.youtube.com/user/jlstine/}
\end{itemize}

\end{document}